\documentclass{article}
\usepackage[utf8]{inputenc}
\usepackage[T1]{fontenc}
\usepackage[margin=2cm]{geometry}
\usepackage{hyperref}
\sloppy
\begin{document}
\section*{Historical source}
Encyclopaedia Britannica scan, source row 611. Parse 3.5 extracted the text. Super 3.5 VL corrected it from an enlarged crop of each block. It can still contain recognition errors.
Page headers, footers, and page numbers are removed.

\subsection*{Untitled continuation (E0060)}
\textit{Continues an entry from an earlier page.}

Refraction. curvature will be inefenfible, and the fluid included in the space ZYHQ will move uniformly as if the solid body were not there. The motions on the other fide of the axis AC will be the fame ; and we need only attend to one half, and we shall confider these as in a state of permanency.

No body changes either its direction or velocity otherwise than by insensible degrees: therefore the particle which is moving in the axis will not reach the vertex A of the body, where it behooved to deflect infantaneously at right angles. It will therefore begin to be deflected at some point F a-head of the body, and will describe a curve FM, touching the axis in F, and the body in M; and then, gliding along the body, will quit it at some point I, describing a tangent curve, which will join the axis again (touching it) in R; and thus there will be a quantity of flagrant water FAM before or a-head of the body, and another LCR behind or a-tern of it.

Let a be the velocity of a particle of the fluid in any instant, and a' its velocity in the next instant. The velocity a may be considered as compounded of a' and a". If the particles tended to move with the velocities a" only, the whole fluid would be in equilibrium (general principle), and the preasure of the fluid would be the same as if all were frangnant, and each particle were urged by a force a''/r, r expressing an indefinitely small moment of time. (N. B. a''/r is the proper expression of the accelerating force, which, by acting during the moment r, would generate the velocity a''; and a'' is supposed an indeterminate quantity, different perhaps for each particle). Now let a be supposed confiant, or a=a'. In this case a''=0. That is to say, no preasure whatever will be exerted on the solid body unless there happen changes in the velocities or directions of the particles.

Let a and a' then be the motions of the particles in two consecutive instants. They would be in equilibrio if urged only by the forces a''. Therefore if y be the point where the particles which describe the curve FM begin to change their velocity, the pressure in D would be equal to the pressure which the fluid contained in the canal y FMD would exert, if each particle were solicited by its force a''. The question is therefore reduced to the finding the curvature in the canal y FMD, and the accelerating forces a'' in its different parts.

It appears, in the first place, that no preface is exerted by any of the particles along the curve FM: for supposing that the particle a (fig. 19.) describes the indefinitely small straight line ab in the first instant, and bc in the second instant; produce ab till bd=ab, and joining dc, the motion ab or bd may be considered as compounded of bc, which the particle really takes in the next instant, and a motion dc which should be destroyed. Draw bi parallel to dc, and ie perpendicular to bc. It is plain that the particle b, solicited by the forces bc, ei (equivalent to dc) should be in equilibrio. This being established, be must be = o, that is, there will be no accelerating or retarding force at b: for if there

be, draw bm (fig. 20.) perpendicular to bF, and the parallel nq infinitely near it. The part bn of the fluid contained in the canal bnqm would fulain some pressure from b towards n, or from n towards b. Therefore since the fluid in this flagrant canal should be in equilibrio, there must also be some action, at least in one of the parts bm, mq, qn, to counterbalance the action on the part bn. But the fluid is flagrant in the space FAM (in consequence of the law of continuity). Therefore there is no force which can act on bm, mq, qn; and the pressure in the canal in the direction bn or nb is nothing, or the force b=o, and the force ie is perpendicular to the canal; and there is therefore no pressure in the canal FM, except what proceeds from the part y F, or from the force ei; which last being perpendicular to the canal, there can be no force exerted on the point M, but what is propagated from the part y F.

The velocity therefore in the canal FM is constant if finite, or infinitely small if variable: for, in the first case, the force bc would be absolutely nothing; and in the second case, it would be an infinitesimal of the second order, and may be considered as nothing in comparison with the velocity, which is of the first order. We shall fee by and by that the last is the real state of the case. Therefore the fluid, before it begins to change its direction in F, begins to change its velocity in some point y-a-head of F, and by the time that it reaches F its velocity is as it were annihilated.

Cor. 1. Therefore the pressure in any point D arises both from the retardations in the part F, and from the particles which are in the canal MD: as they all move along the surface of the body, the force a, destroyed in every particle, is compounded of two others, one in the direction of the surface, and the other perpendicular to it; call these p and q. The point D is pressed perpendicularly to the surface MD; first, by all the forces p in the curve MD; secondly, by the force q acting on the single point D. This may be neglected in comparison of the indefinite number of the others; therefore taking in the arc MD, an infinitely small portion Nm = s, the pressure on D, perpendicular to the surface of the body, will be = ∫p s; and this fluent must be so taken as to be = 0 in the point M.

Cor. 2. Therefore, to find the pressure on D, we must find the force p on any point N. Let u be the velocity of the particle N, in the direction Nm in any instant, and u + u its velocity in the following instant;

we must have p = - . Therefore the whole question is reduced to finding the velocity a in every point N, in the direction Nw.

And this is the sum of a series of propositions which follow, in which the author displays the most accurate and precise conception of the subject, and great address and elegance in his mathematical analysis. He at length brings out an equation which expresses the preface on the body in the most general and unexceptionable manner. We cannot give an abstract, because the train of reasoning is already concise in the extreme: nor can we even exhibit the final equation; for it is conceived in the


\section*{Source record}
\begin{description}
\item[Scan ID] \texttt{britannica:309f765b1ecc83e0a7f7762e44b4a927a7555443:train:611}
\item[Image file] \texttt{page-0611-4a137a2c6b07.jpg}
\item[Image SHA-256] \texttt{4a137a2c6b0772a82850ed78c5cb2c6ffd9751f4d7fc1345b3f8035665a788ea}
\item[Dataset] NationalLibraryOfScotland/encyclopaedia\_britannica\_illustrated at revision \texttt{309f765b1ecc83e0a7f7762e44b4a927a7555443} (CC0-1.0)
\item[Extraction] model \texttt{nvidia/nemotron-3.5-parse}, raw reply \texttt{extract/raw/page-0611-4a137a2c6b07.json}
\item[Raw reply SHA-256] \texttt{597125df72cb1e72c83d5bdb98d3e02bd7fba1da626875f77cfc61d3b96f3e20}
\item[Correction] model \texttt{nvidia/Nemotron-3.5-Super-VL}, block crops, raw replies \texttt{correct/raw/crop\_nothink/page-0611-4a137a2c6b07-b*.json}; super\_vl: 10
\item[Running header] R E S 752 T RES
\end{description}
\end{document}
